\section{SWE-bench：软件工程评测的基准}
\subsection{Lite 与 Full 的演进}
Graham 接下来介绍 SWE-bench：一个对 Agent 执行 GitHub 实际 Issue 的评测框架。每个样本包含原始仓库、Issue 描述与测试套件，Agent 的任务是理解语义和上下文，编写补丁并通过该仓库内的测试以验证正确性。

SWE-bench 分为两类：Lite（300 道相对简单的问题）与 Full（2,294 道来自真实仓库的 Issue）。Lite 集中在理解命令、补丁与基本测试；Full 把问题延展到多文件依赖、跨模块调用与复杂环境。Agent 会被限制在一个上下文窗口内，因此需要设计有效的提示、索引检索与局部推理策略。

\subsection{评估流程与指标}
\begin{figure}[H]
\centering
\includegraphics[width=\textwidth]{frame-02-swebench.jpg}
\caption{SWE-bench 的评测统计板块，展示 Lite 与 Full 数据集的成功率。 \protect\footnotemark}
\end{figure}
\footnotetext{视频画面时间区间：00:15:20--00:15:25。}

评估流程包含三步：1) Agent 阅读 Issue 并在仓库中检索相关上下文；2) 生成并应用补丁（编辑文件、运行 \texttt{git diff}）；3) 运行目标测试，并把结果与输出日志存档。系统按照 \texttt{pass/fail} 统计成功率，并把失败样本回馈训练流程以改进提示与工具策略。

\begin{table}[H]
\centering
\begin{tabular}{p{0.35\textwidth} p{0.58\textwidth}}
\toprule
维度 & 描述 \\
\midrule
任务类型 & 从 Issue 中提取目标、定位相关文件、生成补丁并运行对应测试。 \\
\midrule
Lite & 300 个易错点，旨在验证 prompt 结构与快速修复能力，当前最佳 Agent 通过率在 40\%--50\% 之间。 \\
\midrule
Full & 2,294 个挑战，包括异步调用、依赖迁移与安全补丁，目前整体通过率低于 20\%。 \\
\bottomrule
\end{tabular}
\caption{SWE-bench 的设计维度与当前效果。}
\end{table}

\begin{knowledgebox}{SWE-bench 的评估侧重点}
SWE-bench 强调工程闭环：Agent 必须提交 patch、运行测试并提供日志，不能只给出 diff；Lite/Full 的划分让研究团队能先优化基础能力再扩展到复杂场景。
\end{knowledgebox}

\begin{importantbox}{量化指标}
Graham 强调：Lite 级别的 40\% 通过率已经意味着 Agent 能替代部分重复任务，但 Full 级别的低通过率提示我们仍处在 “半自动化” 时代。
\end{importantbox}

\begin{warningbox}{覆盖率局限}
SWE-bench 主要评估 coding+测试，不包含设计、代码审查、部署和维护等更广的软件工程任务，所评估的能力仍然偏重修复型问题。
\end{warningbox}

\subsection{本章小结}
\begin{itemize}
    \item SWE-bench 提供了从简单到复杂的双层跨度，让 Agent 具备逐步演进的目标。
    \item 当前指标仍然远低于人类工程师，在更广的软件生命周期阶段仍需新增评测。
\end{itemize}

